\documentclass[10pt]{article}
\usepackage[utf8]{inputenc}
\usepackage[T1]{fontenc}
\usepackage{amsmath}
\usepackage{amsfonts}
\usepackage{amssymb}
\usepackage[version=4]{mhchem}
\usepackage{stmaryrd}

\begin{document}
Applications of IC

\section*{Applications of the IC characterization}
\begin{itemize}
  \item Revenue equivalence.
  \item Obtaining equilibria in specific auctions.
  \item Optimal reserve prices and entry fees.
  \item Myerson (1981), ( $N=1, I$ bidders)
  \item Alternative expression for revenue.
  \item Optimal mechanism.
  \item Optimal reserve price in classic auctions.
\end{itemize}

\section*{Revenue equivalence}
Theorem 1 (Revenue equivalence) IPVM with $I$ bidders. $X_{i}=\left[\underline{x}_{i}, \bar{x}_{i}\right]$. Let $(p, t)$ and $\left(p^{\prime}, t^{\prime}\right)$ be two DRMs and let $(Q, T),\left(Q^{\prime}, T^{\prime}\right)$ be the corresponding expected probability of trade and expected transfer functions. If

$$
(\forall i) Q_{i}=Q_{i}^{\prime} \text {, then } T_{i}=T_{i}^{\prime}+\left[U_{i}^{\prime}\left(\underline{x}_{i}\right)-U_{i}\left(\underline{x}_{i}\right)\right] \text {. }
$$

\begin{itemize}
  \item This formulation of revenue equivalence and its proof in the next slide are due to Williams (1999).
  \item The same statement and proof hold in the extension of the model to $N$ goods.
\end{itemize}

Proof. By definition,

$$
\begin{aligned}
T_{i}\left(x_{i}\right) & =\frac{\mathrm{d} U_{i}\left(x_{i}\right)}{\mathrm{d} x_{i}} \cdot x_{i}-U_{i}\left(x_{i}\right) \\
T_{i}^{\prime}\left(x_{i}\right) & =\frac{\mathrm{d} U_{i}^{\prime}\left(x_{i}\right)}{\mathrm{d} x_{i}} \cdot x_{i}-U_{i}^{\prime}\left(x_{i}\right)
\end{aligned}
$$

\begin{itemize}
  \item By BIC, $\frac{\mathrm{d} U_{i}}{\mathrm{~d} x_{i}}=Q_{i}=Q_{i}^{\prime}=\frac{\mathrm{d} U_{i}^{\prime}}{\mathrm{d} x_{i}}$ a.e.
  \item Then, $T_{i}\left(x_{i}\right)=T_{i}^{\prime}\left(x_{i}\right)+\left[U_{i}^{\prime}\left(x_{i}\right)-U_{i}\left(x_{i}\right)\right]$ a.e.
  \item If two convex functions have the same gradient a.e., then both functions are equal up to a constant Rockafellar (1970), Chapter 24. Hence,
\end{itemize}

$$
\left(\forall x_{i}\right)\left[U_{i}^{\prime}\left(x_{i}\right)-U_{i}\left(x_{i}\right)\right] \text { is a constant. }
$$

QED

\section*{Revenue equivalence}
Corollary 1 The four classic auctions yield the same expected revenue to the seller.

Proof. At the respective equilibria of the four classic auctions, every bidder $i$ faces $Q_{i}\left(x_{i}\right)=F\left(x_{i}\right)^{I-1}$ and $U_{i}\left(\underline{x}_{i}\right)=0$. QED

\section*{First-price sealed-bid auction}
Theorem 2 (FPA) Symmetric IPVM, $N=1$ good, $I$ bidders. The unique, symmetric, monotone Bayesian-Nash equilibrium of the FPA is

$$
(\forall i), b_{i}\left(x_{i}\right)=x_{i}-\int_{0}^{x_{i}}\left(\frac{F(y)^{I-1}}{F\left(x_{i}\right)^{I-1}}\right) \mathrm{d} y .
$$

Proof. Sketch.

\begin{enumerate}
  \item If the equilibrium strategy is monotone nondecreasing then
\end{enumerate}

$$
Q_{i}\left(x_{i}\right)=F\left(x_{i}\right)^{I-1}
$$

\begin{enumerate}
  \setcounter{enumi}{1}
  \item Under the assumption of monotonicity, we have the system
\end{enumerate}

$$
\begin{aligned}
U_{i}\left(x_{i}\right) & =F\left(x_{i}\right)^{I-1} x_{i}-T_{i}\left(x_{i}\right) \\
U_{i}\left(x_{i}\right) & =U_{i}(0)+\int_{0}^{x_{i}} F(y)^{I-1} \mathrm{~d} y \\
T_{i}\left(x_{i}\right) & =F^{I-1}\left(x_{i}\right) b\left(x_{i}\right)
\end{aligned}
$$

\begin{enumerate}
  \setcounter{enumi}{2}
  \item Solve for $b\left(x_{i}\right)$ and check it is monotone.
\end{enumerate}

\section*{Reserve price and entry fees}
\begin{itemize}
  \item Before bidding agents pay an entry fee $\boldsymbol{c}$. Only agents that choose to submit bids pay $c$.
  \item Generally $c>0$. Sometimes $c<0$ to compensate bidders for cost of bid preparation and to increase participation.
  \item A reserve price $r$ is a minimum price chosen by the seller. It may be announced or kept secret.\\
If no bid exceeds $r$ the seller keeps the object.
  \item In SPA, if only one bid exceeds $r$, the bidder gets the object and pays $r$.
\end{itemize}

\section*{SPA: Reserve price and entry fee}
\begin{itemize}
  \item SIPVM, $X_{i}=[0,1]$. SPA with entry fee $c$ and reserve price $r$.
  \item Buyers' expected profits from bidding are negative if $c>1-r$.
  \item Suppose $0 \leq c \leq 1-r \leq 1$. Find a non-trivial equilibrium of the SPA.
  \item A participant will only bid if her valuation exceeds a marginal value $x_{i}^{\prime}$. The marginal value $x_{i}^{\prime}$ can be determined intuitively by considering the incentives of a bidder with valuation $x_{i}^{\prime}$. Then, use $x_{i}^{\prime}$ to guess an equilibrium strategy.
\end{itemize}

\section*{References}
Myerson, Roger. 1981. "Optimal Auction Design." Mathematics of Operations Research 6: 58-73.\\
Rockafellar, R. Tyrrell. 1970. Convex Analysis. Number 28 in Princeton Mathematical Series. Princeton University Press.\\
Williams, Steven. 1999. "A Characterization of Efficient, Bayesian Incentive Compatible Mechanisms." Economic Theory 14: 155-80.


\end{document}